% Capítulo 5 ampliado con los códigos de entrenamiento, inferencia, servidor,
% adaptador ROS 2 y reglas CLIPS revisados el 10 de septiembre de 2026.
% Los comentarios TODO identifican información por confirmar; no se imprimen.
% Los ejemplos ilustran contratos y transformaciones, no resultados experimentales.
% Conservar las referencias a los capítulos 3 y 4 al integrar en la tesis.

\chapter{Implementación}
\label{ch:implementacion}

Este capítulo describe los componentes desarrollados para transformar
instrucciones en lenguaje natural en planes de acciones para Justina.
La implementación concreta las decisiones metodológicas del
Capítulo~\ref{ch:metodologia} mediante la generación de datos supervisados,
la especialización del intérprete Qwen, la conversión de objetivos a
hechos simbólicos y la expansión de acciones con CLIPS.

El sistema se organiza en componentes con responsabilidades diferenciadas.
La construcción del conjunto de datos y el entrenamiento se realizan fuera
del ciclo de operación del robot. Durante la planificación, el intérprete
recibe una instrucción y devuelve objetivos que se validan y convierten
antes de ingresar al motor simbólico. El plan resultante se publica para
su consumo por el sistema de ejecución. La arquitectura GPT directo se
trata como una implementación separada, con su propio contrato de
respuesta e interfaz de generación.

Como ejemplo conductor se utiliza la orden de búsqueda, adquisición y
entrega de un objeto introducida en el capítulo anterior. Sus diferentes
representaciones permiten relacionar la entrada lingüística, los objetivos,
los hechos CLIPS y las acciones del ejecutor. Los resultados cuantitativos
de entrenamiento y evaluación se reservan para el capítulo de resultados.

\section{Plataforma experimental}
\label{sec:plataforma_implementacion}

\subsection{Robot Justina}

La implementación se integra en la arquitectura ViRBot de Justina,
descrita en la Sección~\ref{subsec:integracion_virbot}. El sistema de
planificación utiliza las capacidades de navegación, percepción,
manipulación e interacción disponibles mediante las interfaces del robot.
Su salida especifica las acciones que deben solicitarse a esos componentes.

La comunicación con la plataforma se realiza mediante ROS 2. El
planificador produce identificadores de acciones y parámetros simbólicos,
como una ubicación de navegación, un objeto o una referencia a una persona.
La resolución física de esas instrucciones corresponde al ejecutor y a los
módulos encargados de cada tarea.

En el ciclo GPSR se utiliza \texttt{instruction\_point} para representar
el punto de interacción con el operador. Este símbolo se distingue de
\texttt{inspection point}, que puede aparecer como una entidad del
catálogo. La diferencia se conserva al normalizar los parámetros del plan.

% TODO: Añadir únicamente las características físicas de Justina relevantes
% para las pruebas realizadas: sensores, base, manipuladores y computadora
% a bordo. Confirmar la configuración concreta utilizada en el experimento.

\subsection{Plataforma de cómputo}

Para el trabajo con Qwen2.5-7B se utilizó una computadora con una GPU
ASUS RTX 4070 Ti, un procesador Intel Core i7 de undécima generación y
32 GB de memoria RAM. La adaptación del modelo se configuró considerando
los recursos disponibles en este equipo.

La implementación admite dos modalidades de despliegue de la inferencia:
local y remota. En la modalidad local, Qwen se ejecuta en el mismo equipo
que controla al robot y comparte sus recursos de cómputo con los módulos
de navegación, percepción e interacción. Esta modalidad permite realizar
la interpretación sin recurrir a un servidor externo, siempre que el
equipo disponga de memoria y capacidad de procesamiento suficientes
para las tareas concurrentes.

Cuando la carga conjunta exige liberar recursos del equipo del robot,
el modelo puede desplegarse como un servicio en otra computadora con GPU,
accesible mediante una red local o a través de Internet. En esta modalidad,
Justina envía la orden al servidor por HTTP y recibe los objetivos
semánticos generados. El modelo permanece cargado en el servidor entre
solicitudes, mientras que la conversión a hechos, la expansión simbólica
y la comunicación con el ejecutor se realizan en el entorno ROS 2 del
robot.

Ambas modalidades conservan la misma función del intérprete y el mismo
contrato de objetivos. La elección depende de los recursos disponibles
y de la carga de los demás componentes. El despliegue remoto permite
trasladar el cómputo de inferencia a otro equipo e incorpora el tiempo de
comunicación de red al recorrido de la solicitud. Por ello, la modalidad
utilizada y la ubicación de los procesos forman parte de la configuración
experimental al medir latencia y uso de recursos.

% TODO: Confirmar el modelo exacto de GPU y CPU, la capacidad de VRAM y qué
% tareas se ejecutaron en cada equipo. El hardware indicado procede del
% contexto de la tesis; no se han añadido especificaciones no confirmadas.

\subsection{Entorno de software y ROS 2}

Los componentes de generación, validación e interpretación se desarrollaron
en Python. El entrenamiento utiliza PyTorch para el cálculo de tensores,
Transformers para cargar Qwen y gestionar el ajuste mediante
\texttt{Trainer}, PEFT para incorporar los adaptadores LoRA y
bitsandbytes para cargar los pesos cuantizados. La preparación de las
particiones utiliza la biblioteca Datasets.
La interacción con CLIPS se implementó mediante \texttt{clipspy}, cuyo
módulo Python se importa como \texttt{clips}.

ROS 2 proporciona los servicios y tópicos que conectan el intérprete,
el adaptador y el planificador con Justina. Los archivos de lanzamiento
agrupan los nodos y configuran sus nombres, parámetros y espacios de
nombres. Esta organización permite ejecutar pruebas de componentes por
separado y después utilizar el mismo contrato en el sistema integrado.

La reproducibilidad requiere identificar las versiones del sistema
operativo, la distribución de ROS 2, Python y las bibliotecas de
entrenamiento e inferencia. Estas versiones se vinculan con el entorno
efectivamente utilizado, ya que las configuraciones de desarrollo y de
despliegue pueden diferir.

% TODO: Completar versiones de Ubuntu, ROS 2, Python, PyTorch, Transformers,
% PEFT, bitsandbytes, Accelerate, CUDA, controlador NVIDIA y clipspy.
% Usar los registros del entorno del experimento, no versiones instaladas
% posteriormente en otros equipos.

\section{Implementación del conjunto de datos}
\label{sec:implementacion_dataset}

\subsection{Generador de órdenes}

La generación se organiza alrededor de \texttt{generate\_dataset.py},
que coordina la selección de muestras, el balance, la deduplicación y
la validación previa a la escritura. El componente
\texttt{gpsr\_commands.py} construye las órdenes utilizando las familias
y variantes definidas en \texttt{command\_constants.py}.

La generación semántica se implementó en \texttt{command\_goals.py}
mediante la clase \texttt{CommandGoalsMixin}. El método
\texttt{\_generate\_goals} selecciona la función registrada para una
familia y le entrega el contexto de construcción. Cada función devuelve
una lista ordenada de objetivos. Las continuaciones de una orden pueden
invocar otras funciones y recibir el objeto o la persona establecidos
por las subtareas anteriores.

Por ejemplo, la función de transporte entre dos ubicaciones toma el
objeto, el origen y el destino del contexto y construye los objetivos
de navegación, búsqueda, adquisición y colocación. La función de entrega
al operador utiliza el objeto conservado en el contexto y genera
\texttt{deliver(object, to=me)}. Esta propagación permite mantener el
referente de un pronombre en las distintas partes de una orden.

La representación de entrada y etiqueta del ejemplo conductor es:

{\scriptsize
\begin{verbatim}
{
  "input": "Navigate to the bed then find a chocolate jello and take it and give it to me",
  "goals": [
    "go(bed)",
    "find(chocolate jello, kind=object)",
    "take(chocolate jello)",
    "deliver(chocolate jello, to=me)"
  ]
}
\end{verbatim}
}

El ejemplo omite los metadatos para destacar la correspondencia entre
entrada y objetivos. En el archivo JSONL, cada muestra completa se
serializa en una sola línea.

\subsection{Catálogos y CompetitionTemplate}

El módulo \texttt{knowledge.py} carga los catálogos mediante
\texttt{parse\_data}. La estructura de entrada separa los nombres de
personas, las ubicaciones, las habitaciones y los objetos en archivos de
CompetitionTemplate. Las rutas relativas utilizadas son:

\begin{verbatim}
names/names.md
maps/location_names.md
maps/room_names.md
objects/objects.md
\end{verbatim}

El conocimiento cargado se entrega al generador para sustituir los
elementos variables de las plantillas. La validación de catálogos se
encapsula en \texttt{catalog\_validation.py}; incluye la detección de
colisiones y diferencias con el catálogo de nombres del banco de
conocimiento de Justina.

La localización del espacio de trabajo del robot puede configurarse con
\texttt{JUSTINA\_WS}. Esto permite vincular las comprobaciones con el
repositorio utilizado en el despliegue sin incorporar rutas personales
al contrato de generación.

\subsection{Validadores}

La validación de etiquetas se implementó en \texttt{goal\_schema.py}.
La función \texttt{parse\_goal} separa el nombre del objetivo, los
argumentos posicionales y los campos nombrados. El resultado se almacena
en una estructura \texttt{ParsedGoal}, que permite consultar el objetivo
principal y el tipo de entidad sin depender de su capitalización.

La función \texttt{validate\_goals} aplica comprobaciones individuales
y de secuencia. Detecta sintaxis inválida, ausencia de argumentos
obligatorios, campos no admitidos, incompatibilidades entre calificadores
y tipos, y algunas dependencias incumplidas. Las incidencias se
representan mediante \texttt{ValidationIssue}, con un código estable,
un mensaje y, cuando corresponde, la posición del objetivo.

La auditoría del conjunto se implementó en
\texttt{dataset\_evaluation.py}. Su función
\texttt{evaluate\_dataset\_rows} recorre las filas, valida sus
objetivos y acumula estadísticas por familia, firma, tipo y campos.
El módulo también contiene \texttt{ClipsPlanValidator}, que carga las
reglas una vez y reinicia el entorno para cada secuencia examinada.

El programa \texttt{evaluate\_dataset.py} permite ejecutar esa auditoría
sobre un JSONL. Los límites de lectura y de comprobación con CLIPS se
configuran de manera independiente. Cuando se fija un límite de
secuencias CLIPS, se toman las primeras etiquetas válidas en el orden
de lectura hasta alcanzarlo. El reporte agregado se imprime al finalizar;
las etiquetas inválidas o los fallos de expansión producen un código de
salida de error.

Las pruebas de \texttt{test\_dataset\_quality.py} comprueban propiedades
del contrato y del generador sobre un lote de 500 muestras. Incluyen
cardinalidad, deduplicación, distribución entre categorías y firmas
esperadas para familias concretas. La prueba de expansión utiliza el
primer ejemplo de cada familia presente en ese lote. Este tamaño de
prueba se distingue de la cantidad de muestras comprobadas en la
auditoría del conjunto final.

% TODO: Indicar la cobertura y el resultado de la auditoría realmente
% ejecutada. La validación CLIPS confirmada fue sobre algunas muestras;
% no se atribuye al dataset final el tamaño del lote de pruebas unitarias.

\subsection{Dataset final}

El conjunto se almacena en \texttt{dataset\_gpsr.jsonl}, con un objeto
JSON por línea. Cada fila contiene \texttt{input}, \texttt{goals} y
\texttt{meta}. Los objetivos corresponden al vocabulario de la
Tabla~\ref{tab:vocabulario_goals}. La división considerada comprende
36\,000 muestras de entrenamiento y 4\,000 de validación.

Los metadatos identifican la familia, la categoría, la variante
superficial, la firma de objetivos, los tipos de entidades y los campos
presentes. Se conservan para caracterizar el conjunto, pero se excluyen
del contenido supervisado del modelo. No contienen los resultados
individuales de la validación con CLIPS.

El entrenamiento utiliza la instrucción como entrada y un objeto con
la propiedad \texttt{goals} como respuesta esperada. El cargador conserva
temporalmente la familia y la categoría como columnas auxiliares; estas
columnas se eliminan al construir las entradas tokenizadas y no forman
parte de la conversación supervisada.

Después de validar las filas, se aplica \texttt{train\_test\_split}
con una proporción de validación de 0.1 y semilla 42. La partición se
realiza aleatoriamente por filas, sin agrupación ni estratificación por
familia o plantilla. Por ello, la pérdida de validación caracteriza
ejemplos de la distribución del generador y no demuestra, por sí sola,
generalización a familias o plantillas excluidas del entrenamiento.

% TODO: Confirmar los recuentos y la distribución por familia del archivo final.
% Registrar la huella del dataset y su correspondencia con el adaptador.

\section{Entrenamiento e inferencia de Qwen2.5-7B}
\label{sec:implementacion_qwen}

\subsection{Preparación del modelo}

El entrenamiento se organiza en \texttt{Nl-Cl.py} y selecciona el
repositorio \texttt{Qwen/Qwen2.5-7B} como modelo base. Cada fila se valida
y se transforma en una conversación mediante
\texttt{apply\_chat\_template} del tokenizador. El mensaje del usuario
contiene la orden y la respuesta del asistente contiene los objetivos
serializados como JSON. Esta preparación utiliza esos dos roles, sin
añadir un mensaje de sistema al ejemplo supervisado.

El cálculo de la pérdida se concentra en la respuesta. Las posiciones
del prompt reciben la etiqueta de exclusión \texttt{-100}, al igual
que el relleno añadido durante la preparación de los lotes. El
componente \texttt{CompletionOnlyCollator} realiza el relleno dinámico
a la derecha, hasta un múltiplo de ocho posiciones, y ajusta las máscaras
de atención y las etiquetas. Si el tokenizador carece de token de
relleno, se utiliza su token de fin de secuencia.

La conversación completa se tokeniza con truncamiento y longitud máxima
de 512 tokens. El prefijo del usuario se tokeniza por separado para
determinar las posiciones que se excluyen de la pérdida. El límite
comprende tanto la orden como la respuesta, por lo que puede afectar a
los objetivos situados al final de una secuencia extensa.

% TODO: Registrar la revisión exacta del modelo base y del tokenizador.
% Cuantificar ejemplos truncados y verificar que conservan tokens de respuesta
% supervisados; el script no aporta por sí solo ese recuento.

\subsection{Configuración QLoRA}

El modelo base se carga con cuantización de cuatro bits en formato NF4
y doble cuantización. Sus pesos permanecen congelados y el ajuste se
realiza sobre adaptadores LoRA insertados en las proyecciones de atención
y del bloque de transformación del modelo.

Los módulos seleccionados son \texttt{q\_proj}, \texttt{k\_proj},
\texttt{v\_proj}, \texttt{o\_proj}, \texttt{gate\_proj},
\texttt{up\_proj} y \texttt{down\_proj}. Se utiliza rango 32, factor
de escala 64 y una probabilidad de descarte de 0.05. Esta configuración
delimita la capacidad entrenable sin actualizar todos los parámetros
del modelo base.

El cálculo utiliza BF16 cuando la GPU declara soporte para ese formato
y FP16 en caso contrario. La carga selecciona la implementación de
atención \texttt{eager}, desactiva la ventana deslizante y distribuye
el modelo con \texttt{device\_map='auto'}. Durante el ajuste se
desactiva la caché y se habilita \textit{gradient checkpointing} con
\texttt{use\_reentrant=False}, reduciendo las activaciones conservadas
a cambio de recalcularlas durante la retropropagación. El script requiere
CUDA y limita a 0.9 la fracción de memoria disponible para el asignador
de PyTorch en el dispositivo configurado.

Al guardar el modelo se conserva el adaptador y su configuración,
junto con el tokenizador. La inferencia requiere cargar estos artefactos
sobre la misma versión del modelo base.

\subsection{Hiperparámetros de entrenamiento}

La Tabla~\ref{tab:hiperparametros_qwen} reúne la configuración de
entrenamiento establecida en el script de especialización. La
acumulación de gradientes permite procesar una muestra por GPU y
actualizar los parámetros después de acumular ocho muestras.

\begin{table}[htbp]
\centering
\caption{Configuración de entrenamiento de Qwen2.5-7B.}
\label{tab:hiperparametros_qwen}
\begin{tabular}{|p{0.58\textwidth}|p{0.28\textwidth}|}
\hline
\textbf{Parámetro} & \textbf{Valor} \\
\hline
Longitud máxima de entrada y respuesta & 512 tokens \\
\hline
Tamaño de lote por GPU & 1 \\
\hline
Acumulación de gradientes & 8 \\
\hline
Tamaño efectivo de lote, con una GPU & 8 \\
\hline
Número de épocas & 2 \\
\hline
Tasa de aprendizaje inicial & \(1\times10^{-4}\) \\
\hline
Planificador de tasa de aprendizaje & Coseno \\
\hline
Pasos de calentamiento & 150 \\
\hline
Intervalo de evaluación & 250 pasos \\
\hline
Intervalo de guardado & 500 pasos \\
\hline
Intervalo de registro de pérdida & 50 pasos \\
\hline
Rango / escala / descarte LoRA & 32 / 64 / 0.05 \\
\hline
\end{tabular}
\end{table}

Con 36\,000 muestras de entrenamiento, un lote efectivo de ocho y dos
épocas, la configuración corresponde a 9\,000 actualizaciones. Este
valor es una consecuencia de la configuración y no sustituye el registro
de la ejecución efectivamente realizada.

El optimizador y la semilla del entrenamiento no se sobrescriben
explícitamente en \texttt{TrainingArguments}; su valor efectivo depende
de la configuración resuelta por la versión de Transformers utilizada.
La semilla 42 de la partición identifica una operación distinta de las
fuentes de aleatoriedad del ajuste.

% TODO: Contrastar la configuración con training_args.bin o su exportación
% y trainer_state.json de la ejecución asociada al dataset definitivo.
% Registrar optimizador, precisión efectiva y semillas resueltas.

\subsection{Selección del checkpoint}

La selección automática se configura mediante estas opciones del
entrenador:

\begin{verbatim}
load_best_model_at_end=True
metric_for_best_model="eval_loss"
greater_is_better=False
\end{verbatim}

Al terminar, el entrenador restaura
el checkpoint seleccionado por menor pérdida de validación entre los
candidatos que registra para ese propósito. La evaluación se programa
cada 250 actualizaciones y el guardado cada 500; estos intervalos deben
distinguirse al relacionar una pérdida con un archivo conservado. Se
configura además \texttt{save\_total\_limit=2} para limitar la retención.

Después de \texttt{trainer.train()}, se guardan el adaptador seleccionado
y el tokenizador en \texttt{nl2cd\_qwen7b}. La curva de entrenamiento
se construye a partir de \texttt{trainer.state.log\_history}, utilizando
los registros de pérdida de entrenamiento y validación, y se exporta
como \texttt{loss\_curve.png}.

El script incluye también una comprobación sobre las primeras 50
muestras de validación, o sobre las disponibles si fueran menos. Este
límite pertenece al código de comprobación posterior al ajuste y no
define el conjunto experimental de comparación. La pérdida de
validación, la coincidencia de objetivos y la expansión con CLIPS
aportan evidencias distintas y se interpretan por separado.

El artefacto destinado a inferencia se identifica de manera conjunta
por el modelo base, el adaptador y el tokenizador. El nombre de un
directorio de salida no basta para determinar qué checkpoint contiene;
la selección debe quedar vinculada con su paso y configuración de
entrenamiento.

% TODO: Añadir best_model_checkpoint, best_metric y global_step del registro
% efectivo, y comprobar qué adaptador se desplegó. El código determina el
% procedimiento configurado, no la identidad de una ejecución histórica.

\subsection{Inferencia}

El módulo \texttt{inference.py} carga \texttt{Qwen/Qwen2.5-7B} en NF4
con doble cuantización y añade los pesos mediante
\texttt{PeftModel.from\_pretrained}. El tokenizador se recupera del
directorio del adaptador y el modelo se coloca en modo de evaluación.
El dispositivo predeterminado es \texttt{cuda:0}; las opciones de
ejecución permiten cambiar la distribución del modelo y establecer
límites de memoria.

Antes de tokenizar, el recorrido de inferencia aplica el normalizador
obtenido mediante \texttt{get\_default\_normalizer} del módulo
\texttt{command\_normalizer}. La orden resultante se introduce como
mensaje de usuario y la plantilla añade el inicio de la respuesta del
asistente. La entrada normalizada se devuelve junto con la predicción
cuando se solicita \texttt{return\_normalized=True}.

La generación utiliza \texttt{do\_sample=False}, penalización de
repetición de 1.05 y un máximo de 128 tokens nuevos. Las condiciones de
parada incluyen el token de fin de secuencia y \texttt{<|im\_end|>},
cuando está disponible. El script de entrenamiento utiliza un límite
de 256 tokens en sus comprobaciones posteriores al ajuste; por ello,
ambos recorridos no deben tratarse como configuraciones idénticas de
inferencia.

Se retiran los tokens de entrada de la secuencia producida y se
decodifica únicamente la continuación. La extracción localiza el primer
fragmento con llaves balanceadas e intenta convertirlo mediante
\texttt{json.loads}. Es un procedimiento de extracción seguido de
validación, sin restricción gramatical durante la generación. Si la
extracción o la validación falla, se devuelve un error y hasta 200
caracteres de la respuesta como información diagnóstica.

La salida se procesa para obtener el objeto JSON y se somete a la
validación del contrato semántico. El resultado aceptado contiene una
lista \texttt{goals}; las respuestas que incumplen el contrato se
rechazan antes de entregarlas al siguiente componente. Esta comprobación
detecta errores de estructura y restricciones programadas, sin acreditar
por sí sola la fidelidad de la interpretación.

La predicción directa corresponde a
\(G_{\mathrm{pred}}^{\mathrm{Qwen}}\). Las normalizaciones o reparaciones
posteriores se consideran operaciones adicionales del sistema y se
distinguen del resultado inicial al evaluar el intérprete.

% TODO: Incorporar command_normalizer.py para especificar sus transformaciones.
% Fijar el límite de generación utilizado en la comparación y conservar la
% salida original cuando se requiera un diagnóstico más allá de 200 caracteres.

\subsection{Despliegue del modelo}

El intérprete Qwen admite dos modalidades de uso: ejecución local y
ejecución como servicio remoto. Ambas utilizan el modelo base, el
adaptador LoRA y el tokenizador, y producen objetivos mediante el mismo
procedimiento de inferencia. La elección determina dónde se realiza
el cómputo y cómo se comunica la orden al intérprete.

En la modalidad local, el modelo se carga en el mismo equipo que
ejecuta los componentes de control del robot. El módulo
\texttt{inference.py} proporciona las funciones de carga y traducción,
por lo que la interpretación puede realizarse directamente desde un
proceso Python. Una vez disponibles los archivos del modelo y del
adaptador, el procesamiento de cada orden no requiere consultar un
servidor externo. Esta modalidad comparte los recursos de CPU, GPU y
memoria con los demás componentes de Justina.

Cuando la carga de trabajo requiere liberar recursos del equipo del
robot, la inferencia puede trasladarse a otra computadora y exponerse
como un servicio accesible mediante una red local o Internet. Para
esta modalidad se desarrolló \texttt{server.py}, que mantiene cargados
el modelo, el tokenizador y el normalizador entre solicitudes. Justina
envía la instrucción por HTTP y recibe la predicción, mientras que
la conversión a hechos y la planificación con CLIPS permanecen en
el entorno ROS 2 del robot.

El servidor utiliza \texttt{ThreadingHTTPServer} y escucha por defecto
en el puerto 8008. Aunque las solicitudes se atienden mediante hilos,
un mecanismo de exclusión mutua permite una sola traducción a la vez.
Esto evita que varias solicitudes ejecuten simultáneamente la
generación sobre el mismo modelo.

La interfaz ofrece dos rutas: \texttt{GET /health}, para consultar
la disponibilidad del servicio, y \texttt{POST /translate}, para
solicitar la interpretación de una orden. La petición de traducción
contiene un objeto JSON con la propiedad \texttt{command}:

\begin{verbatim}
{"command": "go to the kitchen"}
\end{verbatim}

Cuando la predicción es aceptada, la respuesta incluye
\texttt{ok=true}, el tiempo \texttt{elapsed\_ms} y un objeto
\texttt{result}. Dentro de este objeto, \texttt{normalized\_input}
conserva la entrada procesada y \texttt{prediction} contiene la lista
\texttt{goals}. La envoltura HTTP se retira antes de entregar al
adaptador el contrato semántico:

\begin{verbatim}
{"goals": ["go(kitchen)"]}
\end{verbatim}

El servidor limita el cuerpo de la petición a 8192 bytes por defecto.
Una traducción aceptada devuelve el código 200; una petición incorrecta
o una orden vacía, 400; una predicción inválida, 422; y un fallo interno,
500. Cuando se configura una clave de acceso, el servidor comprueba
el encabezado \texttt{Authorization} mediante el esquema
\texttt{Bearer} y devuelve 401 si la clave no coincide.

La modalidad de despliegue también determina el alcance de las
mediciones temporales. En el servicio remoto, \texttt{elapsed\_ms}
incluye la espera por el acceso al modelo, la normalización, la
inferencia y el procesamiento de la salida. Excluye la carga inicial
del modelo y el tiempo completo de comunicación observado por el
cliente. Por ello, la evaluación distingue el tiempo de interpretación
del tiempo total de la solicitud y registra si el modelo se ejecutó
localmente o mediante el servicio remoto.

% TODO: Confirmar si la inferencia del experimento fue local, por red
% local o mediante un túnel remoto. No incluir claves o direcciones
% privadas innecesarias en el documento de tesis.

\section{Implementación del planificador CLIPS}
\label{sec:implementacion_clips}

\subsection{Conversión de objetivos a hechos gpsr-goal}
\label{subsec:adaptador_objetivos_impl}

El adaptador \texttt{goals\_to\_clips\_node.py} recibe la lista de
objetivos y construye hechos \texttt{gpsr-goal}. El componente
\texttt{GoalParser} acepta un objeto con la propiedad \texttt{goals}
o directamente una lista JSON. Analiza cada expresión y separa los
argumentos posicionales de los pares \texttt{clave=valor}, respetando
las comas delimitadoras y los valores entre comillas. Cada hecho conserva
la posición, el nombre del objetivo, el objetivo principal y los
campos utilizados para su expansión. La estructura compartida es:

\begin{verbatim}
(deftemplate gpsr-goal
  (slot step (type INTEGER))
  (slot name)
  (slot target (default ""))
  (slot kind (default unknown))
  (slot destination (default ""))
  (slot relation (default ""))
  (slot value (default ""))
  (slot qualifier (default "")))
\end{verbatim}

Para el segundo objetivo del ejemplo conductor, la representación es:

\begin{verbatim}
(gpsr-goal
  (step 2)
  (name find)
  (target "chocolate_jello")
  (kind object)
  (destination "")
  (relation "")
  (value "")
  (qualifier ""))
\end{verbatim}

Los espacios de los argumentos se convierten en guiones bajos para
construir los identificadores de ejecución. Los valores de cadena se
escapan antes de insertarse en CLIPS, mientras que el nombre y el tipo
se representan como símbolos.

La clasificación de entidades utiliza la secuencia completa. Los
destinatarios de entrega, los participantes de guía, seguimiento y saludo,
los atributos visuales y la capitalización de los nombres aportan indicios
para interpretar una búsqueda como búsqueda de persona. El campo
\texttt{kind} del hecho se calcula mediante estas reglas de contexto;
no se copia directamente del argumento homónimo de la predicción.
Por ejemplo, \texttt{go} y \texttt{place} se clasifican como ubicación,
\texttt{take} y \texttt{deliver} como objeto, y las operaciones de
información como \texttt{info}. Esta clasificación cumple una función
operativa en el adaptador y no redefine el vocabulario semántico.

\texttt{ClipsGoalFactBuilder} conserva los destinos indicados por
\texttt{to}, \texttt{at}, \texttt{on} o \texttt{in}, en ese orden
de prioridad. También admite el segundo argumento posicional para
entrega, guía, seguimiento y colocación. La relación principal selecciona
el primer campo disponible entre destino, gesto, postura, vestimenta y
propiedad. El campo \texttt{qualifier} conserva el primer atributo
visual, con prioridad de gesto, postura, vestimenta y propiedad, en la
forma \texttt{clave:valor}. La representación no codifica una conjunción
de varios calificadores visuales.

Además de la normalización léxica, el adaptador aplica una reparación
específica cuando recibe una navegación cuyo argumento corresponde a
una persona. Ese objetivo se transforma en búsqueda. Si inmediatamente
después aparece un seguimiento de la misma persona con destino, se
adelanta el destino como navegación, se añade la búsqueda y se conserva
el seguimiento sin destino. Esta operación modifica la interpretación
de la secuencia y debe registrarse como posprocesamiento, sin atribuir
la corrección al modelo.

Para mantener la notación del capítulo anterior, la salida del intérprete
conserva el símbolo \(G_{\mathrm{pred}}^{\mathrm{Qwen}}\). Se denomina
\(\operatorname{Adaptar}\) a la transformación del adaptador, que comprende el
análisis, las reparaciones, la clasificación y la serialización. El
recorrido implementado queda expresado como:
\begin{equation}
P^{\text{Qwen--CLIPS}} =
\operatorname{CLIPS}\!\left(
\operatorname{Adaptar}\!\left(G_{\mathrm{pred}}^{\mathrm{Qwen}}\right)\right).
\label{eq:recorrido_adaptador_impl}
\end{equation}
La calidad de la interpretación se calcula sobre la predicción anterior
a esta transformación. La evaluación del sistema completo considera
también los hechos y el plan que resultan de ella.

La auditoría fuera de línea utiliza
\texttt{goal\_to\_clips\_fact} para construir el mismo tipo de hecho,
pero emplea el análisis del esquema semántico y no reproduce todas las
heurísticas y reparaciones del nodo ROS 2. En consecuencia, compartir
la plantilla \texttt{gpsr-goal} no garantiza equivalencia entre ambos
recorridos. La comprobación integrada debe conservar la transformación
que efectivamente recibió CLIPS.

% TODO: Fijar las versiones del esquema y del adaptador para el vocabulario
% final. Verificar la equivalencia de conversiones sobre el conjunto de prueba
% o registrar sus diferencias. Las reparaciones pueden cambiar el significado
% y no deben aceptarse como corrección semántica sin una referencia independiente.

\subsection{Representación del estado simbólico}

La planificación comienza con un estado que representa ubicación
desconocida y ausencia de un objeto sostenido:

\begin{verbatim}
(current-state (location "unknown") (holding "nothing"))
\end{verbatim}

Las reglas actualizan este estado al añadir acciones al plan. Los
valores representan efectos previstos, no confirmaciones físicas. Por
ejemplo, la incorporación de una adquisición modifica el objeto
esperado en posesión del robot. La colocación exige que coincida con
el objeto solicitado; las reglas de entrega no imponen esa misma
igualdad como precondición general. La validación de la secuencia de
objetivos y las condiciones de expansión son, por tanto, controles
complementarios.

Otros hechos conservan referencias a personas, el último contexto de
interacción, análisis visuales y el estado del foco. Las referencias
\texttt{guest\_N} y \texttt{customer\_N} permiten reutilizar una entidad
durante distintas subtareas. Los objetos se identifican mediante nombres
normalizados e índices, como \texttt{chocolate\_jello\_1}.

La referencia \texttt{guest\_last} se utiliza para una persona capturada
por \texttt{find\_person}, mientras que \texttt{customer\_N} identifica
una búsqueda por atributo realizada con \texttt{find\_pose}. Una persona
nombrada puede requerir búsqueda anónima, solicitud del nombre y escucha
mediante \texttt{guest\_name\_last}. La asociación efectiva de estos
identificadores con detecciones y respuestas corresponde al ejecutor.

El modo de manipulación se representa mediante
\texttt{manipulation-mode}. Su valor predeterminado está habilitado y
puede cambiarse mediante \texttt{set-manipulation-enabled}. Cuando se
deshabilita, una adquisición produce una explicación hablada y un hecho
\texttt{manipulation-unavailable}; las operaciones dependientes generan
el aviso correspondiente. Esta salida informa una capacidad no disponible
y no acredita la realización de la tarea de manipulación.

\subsection{Reglas de expansión}

Las reglas se agrupan en \texttt{goals\_planning.clp}. Una regla de
procesamiento consume un objetivo únicamente cuando no existe otro
pendiente con un paso menor. Después de expandirlo, retira el hecho
correspondiente y añade las acciones de la subtarea.

La navegación a una ubicación distinta de la registrada produce una
solicitud \texttt{go\_to} y una confirmación hablada. Si ambas ubicaciones
coinciden, el objetivo se consume sin repetir el desplazamiento. La
regla siguiente muestra esta decisión y la actualización del estado:

{\scriptsize
\begin{verbatim}
(defrule process-go
  (declare (salience 200))
  (planning-ready)
  ?g <- (gpsr-goal (step ?step) (name go) (target ?dest))
  (not (gpsr-goal (step ?prev&:(< ?prev ?step))))
  ?st <- (current-state (location ?loc) (holding ?held))
  =>
  (retract ?g)
  (if (neq ?loc ?dest) then
    (append-say (str-cat "I will navigate to " ?dest))
    (append-action go_to ?dest)
    (append-action say (str-cat "I have arrived at " ?dest))
    (retract ?st)
    (assert (current-state (location ?dest) (holding ?held)))))
\end{verbatim}
}

La búsqueda de objetos genera \texttt{find\_object} o una
operación de análisis cuando se utiliza una propiedad admitida. Las
búsquedas de personas recurren a \texttt{find\_person} o
\texttt{find\_pose} y conservan el contexto necesario para la
interacción posterior.

Con la manipulación habilitada, la adquisición utiliza \texttt{take};
la entrega incorpora aproximación y \texttt{give}; la colocación puede
incluir navegación al destino, \texttt{find\_space} y \texttt{drop}.
La selección concreta depende de los argumentos y del estado simbólico.
Las acciones de diálogo y de control del foco complementan estas
operaciones.

Para la colocación se distinguen puntos de navegación y superficies
genéricas. Un destino como \texttt{storage\_rack} puede añadir navegación,
mientras que \texttt{table}, \texttt{shelf} o \texttt{cabinet} se tratan
como soportes en el lugar actual. La función siguiente incorpora la
búsqueda de espacio y la liberación del objeto; el suelo constituye una
excepción que utiliza directamente el argumento \texttt{floor}:

{\small
\begin{verbatim}
(deffunction append-drop-sequence (?obj ?destination)
  (bind ?surface (drop-surface ?destination))
  (if (eq ?surface "floor") then
    (append-action drop "floor")
    else
    (append-action find_space ?surface)
    (append-action drop (object-id ?obj))))
\end{verbatim}
}

En este fragmento, \texttt{drop} identifica la acción del ejecutor que
materializa el objetivo semántico \texttt{place}. La búsqueda de espacio
selecciona una clase de soporte, como mesa, estante o contenedor, y la
liberación utiliza el identificador del objeto.

La guía genera una instrucción hablada, navegación al destino y cierre
de la interacción. El seguimiento recupera o busca a la persona,
incorpora aproximación cuando corresponde y emite mediante \texttt{say}
las claves \texttt{instruct\_follow\_only} o
\texttt{instruct\_follow\_to\_<destino>}. Esta regla no emite una
acción explícita de seguimiento ni actualiza la ubicación al destino.
Su expansión representa el protocolo de interacción; la realización
física del seguimiento depende del tratamiento de esas claves en el
ejecutor y debe verificarse por separado.

Las propiedades de tamaño se traducen a consultas de análisis de
objetos grandes o pequeños. La implementación también utiliza el
tamaño como aproximación para las solicitudes de mayor o menor peso
y menor grosor. Esta correspondencia permite producir una consulta
visual, pero no constituye una medición de peso o grosor ni demuestra
la satisfacción de esas propiedades. Estos casos requieren una
interpretación diferenciada en la evaluación.

El conteo genera \texttt{analyze\_objects}, incorpora retorno al punto
de instrucción cuando corresponde y emite la respuesta
\texttt{count\_category}. Para guardar un nombre se solicita y escucha
\texttt{guest\_name\_last}; para atributos de persona se utilizan
operaciones de análisis o búsqueda por postura. La información
adquirida se comunica mediante las claves de respuesta de \texttt{tell}.
La respuesta a preguntas combina \texttt{listen} y \texttt{say} con
la clave \texttt{answer\_question}, y el saludo utiliza la persona
referenciada y la clave \texttt{greet}.

% TODO: Confirmar el tratamiento de instruct_follow_* en el ejecutor y
% delimitar los casos de peso/grosor en la comparación. Verificar la actualización
% de location tras place cuando el destino representa un soporte genérico.

\subsection{Generación de acciones}

Las reglas construyen hechos intermedios \texttt{action-step}, que
posteriormente se convierten en \texttt{ros-message}. Los anuncios
generales utilizan numeración desde 1 y las acciones operativas desde
1001. Al ordenar la salida, se presenta primero el anuncio de las tareas
y después las operaciones del plan.

En el ejemplo conductor, la secuencia operativa central puede
representarse como sigue. Se omiten aquí los anuncios y confirmaciones
para mostrar la relación con los cuatro objetivos:

\begin{verbatim}
go_to       bed
find_object chocolate_jello_1
take        chocolate_jello_1
go_to       instruction_point
approach    host_1
give        chocolate_jello_1
\end{verbatim}

El objetivo de entrega al operador se expande en retorno al punto de
instrucción, aproximación y entrega física. Esta secuencia es una
ilustración de la expansión; el plan completo incluye también las
acciones de comunicación y finalización establecidas por las reglas.

\subsection{Condición de finalización}

Cuando no quedan objetivos pendientes, las reglas desactivan el foco
si continúa activo, añaden el retorno a \texttt{instruction\_point}
cuando corresponde y generan el mensaje de tarea completada. La
construcción termina con la nota \texttt{GPSR\_DONE}.

La finalización se distribuye en tres reglas. Primero,
\texttt{add-final-return} agrega el cierre cuando no quedan objetivos.
Después, \texttt{serialize-action-step} transforma las acciones en
mensajes únicamente cuando existe \texttt{final-steps-added}. Finalmente,
\texttt{goals-finish} emite la nota de cierre cuando tampoco quedan
hechos \texttt{action-step} por serializar. Las prioridades de estas
reglas son 20, 5 y 1, respectivamente.

El componente \texttt{PlanExtractor} ordena los hechos
\texttt{ros-message} por su paso y construye una cadena con una
instrucción por línea. Las notas \texttt{GPSR\_START} y
\texttt{GPSR\_DONE} delimitan el plan publicado. La salida se añade
también al archivo \texttt{plan\_execution.log}.

Estas notas corresponden a la generación del plan. Su presencia no
constituye una confirmación de ejecución del robot. La comprobación
funcional revisa asimismo que no queden objetivos pendientes y que no
aparezca la respuesta de objetivo no soportado.

La regla de objetivo no soportado consume el hecho y añade una
explicación hablada. Por esta razón, la ausencia de objetivos pendientes
y la marca de cierre pueden coexistir con una incapacidad de expansión.
El diagnóstico necesita examinar también el contenido del plan, incluidas
las respuestas de manipulación no disponible.

\section{Integración Qwen--CLIPS con Justina}
\label{sec:integracion_implementada}

\subsection{Componentes ROS 2}

La integración utiliza un nodo semántico, un adaptador de objetivos y
el nodo que administra CLIPS. El nodo \texttt{qwen\_semantic\_node.py}
actúa como cliente del servidor de inferencia y entrega la representación
con \texttt{goals}. El adaptador coordina la inserción de hechos y el
inicio de la planificación.

El nodo \texttt{clips\_node.py} expone los servicios del motor y publica
el plan. El acceso al entorno se serializa mediante una cola de tareas
procesada por un único trabajador, evitando que los callbacks manipulen
concurrentemente el mismo entorno CLIPS. El componente
\texttt{clips\_engine.py} encapsula la interacción con
\texttt{clips.Environment}.

El lanzamiento de integración agrupa los componentes bajo el espacio
de nombres \texttt{/planning}. La selección del modo de objetivos
utiliza el parámetro \texttt{plan\_name=goals}, que permite cargar
\texttt{goals\_planning.clp}. El formato de respuesta del nodo
semántico se configura como \texttt{goals} para entregar al adaptador
el contrato esperado.

El adaptador utiliza un nodo auxiliar de clientes y un
\texttt{SingleThreadedExecutor} en un hilo dedicado. Mientras el callback
del servicio de conversión espera la inserción de hechos o el inicio de
planificación, ese ejecutor procesa las respuestas de los servicios
dependientes. Así se evita que ambas operaciones dependan del mismo
callback bloqueado.

\subsection{Flujo de comunicación}

La orden ingresa mediante el servicio del parser semántico. El nodo
envía la solicitud al servidor Qwen y recibe la predicción validada.
La lista de objetivos se entrega después al servicio de conversión,
que coordina dos operaciones: reiniciar e insertar el lote de hechos,
y solicitar el inicio de planificación.

En el despliegue bajo \texttt{/planning}, las interfaces principales
son las siguientes:

{\small
\begin{verbatim}
/planning/semantic_parser/semantic_parser
/planning/goals_to_clips/convert
/planning/clips/assert_facts_batch
/planning/clips/start_planning
/planning/clips/results
\end{verbatim}
}

El servicio de conversión utiliza
\texttt{clips\_interfaces/srv/String}: recibe el JSON en
\texttt{request.data} y devuelve el estado en \texttt{response.data}.
La inserción por lote utiliza el tipo
\texttt{clips\_interfaces/srv/StringArray}. El inicio de planificación
se solicita mediante \texttt{std\_srvs/srv/Trigger}. El hecho
\texttt{(start (name action-planning))} activa la inicialización de las
reglas y la expansión de los objetivos cargados.

El plan completo se publica como un mensaje
\texttt{std\_msgs/msg/String}. La aceptación de la solicitud de
conversión indica que se completó la coordinación de servicios; la
obtención y el contenido del plan se verifican sobre el resultado
publicado. El script \texttt{qwen\_to\_plan.py} permite recorrer este
flujo con una sola orden y mostrar la respuesta semántica, el estado
de conversión y el plan.

El adaptador espera hasta tres segundos por la disponibilidad de cada
servicio dependiente y hasta cinco segundos por su respuesta. Comprueba
\texttt{successful} tras la inserción y \texttt{success} tras el inicio
de planificación. Solo devuelve \texttt{OK} si ambas operaciones
responden favorablemente; en caso de fallo devuelve una descripción
de error. Una lista vacía se rechaza antes del envío. Estos límites
corresponden a cada operación y no equivalen a un tiempo máximo global
del recorrido completo.

Los nombres de conversión, inserción e inicio se declaran como rutas
absolutas bajo \texttt{/planning}. El adaptador solicita el reinicio y
la inserción mediante el servicio por lote; el manejo efectivo de la
memoria de CLIPS pertenece al nodo que atiende ese servicio.

% TODO: Confirmar el launch efectivo, remapeos y tipo del servicio semántico
% con qwen_semantic_node.py, y el contrato del lote con clips_node.py.

\subsection{Comunicación con el ejecutor}

El componente \texttt{plan\_manager\_node} recibe el plan, extrae las
líneas que representan pasos y conserva las instrucciones ordenadas.
Cada solicitud a \texttt{/request\_instruction} obtiene la siguiente
instrucción y la publica en \texttt{/instruction\_msg}. Cuando no
existe un plan asignado, la respuesta indica \texttt{NoPlanAssigned}.

La suscripción del gestor utiliza \texttt{/clips/results}; cuando el
planificador publica bajo \texttt{/planning}, el lanzamiento debe
resolver esta correspondencia mediante el remapeo pertinente. La
configuración de nombres forma parte de la integración y permite que
el consumidor reciba el mismo contrato de salida.

La entrega secuencial separa la construcción del plan de su consumo.
La confirmación de éxito de una acción, el tratamiento de fallos y la
decisión de continuar corresponden a las interfaces del sistema de
ejecución. La publicación de la siguiente instrucción, por sí sola,
no acredita que la anterior se haya realizado físicamente.

Durante el desarrollo se utilizaron 48 comandos como prueba de
funcionamiento. Esta comprobación se considera parte de la verificación
del sistema y se distingue del conjunto experimental del capítulo 6.

% TODO: Confirmar si el recorrido efectivo del experimento utiliza
% plan_manager_node o un supervisor distinto. Completar el contrato con
% el ejecutor, las confirmaciones y la respuesta ante errores. No se
% supone un protocolo de realimentación que aún no se haya identificado.

\section{Implementación del planificador GPT}
\label{sec:implementacion_gpt}

La preparación del componente directo comprende la interfaz de solicitud,
el prompt, el contrato de salida y las comprobaciones que anteceden al
consumo del plan. Estos elementos concretan el diseño de la
Sección~\ref{sec:diseno_gpt_directo} y mantienen separados los objetivos
declarados y las acciones generadas.

% TODO: Confirmar el estado de implementación y aportar el cliente real,
% el prompt y el esquema definitivos. Las subsecciones siguientes describen
% el contrato y la integración prevista, sin atribuirles ejecuciones o
% resultados que todavía no se han confirmado.

\subsection{Interfaz con la API}

La interfaz se especifica mediante una solicitud a la API de Responses.
El mensaje de sistema contiene las instrucciones del planificador y
el mensaje de usuario contiene la orden. La solicitud incorpora también
el esquema de respuesta para delimitar el formato de los objetivos y
del plan.

La configuración separa el identificador del modelo del contenido del
prompt y del esquema. Para la evaluación se fijarán estas versiones y
los parámetros de generación. La respuesta deberá asociarse con la
solicitud que la originó, conservando tanto su contenido como la
información de uso disponible y el tiempo observado por el cliente.

El procesamiento distinguirá las respuestas del modelo de los errores
de servicio, interrupciones y límites de espera. Estos últimos se
registrarán como incidencias de la solicitud y no como decisiones de
aclaración o rechazo de la tarea.

% TODO: Incorporar el cliente de API efectivo, modelo, versión del SDK,
% parámetros admitidos por ese modelo, timeout y política de reintentos.

\subsection{System prompt}

Se elaboró un prompt en inglés con el dominio de entidades, las formas
canónicas de los objetivos y las convenciones de expansión. El
contenido especifica cómo resolver referencias y cómo distinguir
personas, objetos, ubicaciones e información. También establece las
dependencias entre búsqueda, adquisición, transporte e interacción.

Las instrucciones operativas definen el estado esperado durante la
generación: ubicación, objeto sostenido, foco de interacción y
referencias a personas. Se incluyen convenciones de identificación,
numeración de pasos y finalización del plan. El vocabulario de salida
se ajusta al contrato de la investigación.

Para evitar que el contexto cambie inadvertidamente entre solicitudes,
el prompt se conserva como un artefacto separado de la orden. Su
versión deberá permanecer fija durante la comparación. Las instrucciones
operativas permiten examinar la generación de GPT bajo las mismas
convenciones de tarea que utiliza el expansor simbólico.

\subsection{JSON Schema y Structured Outputs}

El contrato prevé utilizar JSON Schema mediante Structured Outputs
para restringir la estructura de la respuesta. Los campos principales
se resumen en la Tabla~\ref{tab:contrato_gpt_impl}. El plan contiene
pasos con un número de orden, el identificador del robot, el nombre
de la acción y sus parámetros.

\begin{table}[htbp]
\centering
\caption{Campos del contrato de respuesta del planificador directo.}
\label{tab:contrato_gpt_impl}
\begin{tabular}{|p{0.24\textwidth}|p{0.63\textwidth}|}
\hline
\textbf{Campo} & \textbf{Función} \\
\hline
\texttt{status} & Tipo de respuesta a la instrucción. \\
\hline
\texttt{message} & Pregunta de aclaración o explicación del rechazo. \\
\hline
\texttt{goals} & Secuencia de objetivos semánticos. \\
\hline
\texttt{plan} & Pasos de acciones compatibles con el contrato del ejecutor. \\
\hline
\texttt{start\_note} & Marca de inicio del plan ejecutable. \\
\hline
\texttt{end\_note} & Marca de finalización de la generación. \\
\hline
\end{tabular}
\end{table}

El esquema y los criterios de aceptación cumplen funciones distintas.
La estructura permite verificar campos, tipos y nombres de acciones;
las relaciones entre estado, objetivos y pasos requieren comprobaciones
adicionales. Un objeto JSON bien formado puede contener un destinatario
incorrecto o una secuencia que incumpla las dependencias de la tarea.

% TODO: Revisar response_schema.json para describir sus restricciones
% exactas. No se presupone que el esquema implemente por sí solo todas
% las condiciones semánticas o dependientes de status del prompt.

\subsection{Generación de objetivos y planes}

Para una instrucción ejecutable, el contrato solicita que el modelo
produzca \(G_{\mathrm{pred}}^{\mathrm{GPT}}\) y
\(P^{\mathrm{GPT}}\) en la misma respuesta. Los objetivos permiten
comparar la interpretación con \(G_{\mathrm{ref}}\); los pasos permiten
examinar la realización de la tarea y su compatibilidad con Justina.

El ejemplo conductor requiere identificar el objeto y su ubicación,
adquirirlo y entregarlo al operador. La especificación directa debe
conducir a acciones con los mismos tipos de identificadores utilizados
en la Sección~\ref{sec:implementacion_clips}, aunque la producción del
plan ocurra dentro del modelo.

La salida de GPT no se introduce en CLIPS. Su consumo requiere validar
el contrato y transformar la representación de los pasos al formato
que recibe el ejecutor. Esa transformación de formato debe distinguirse
de una modificación del contenido operativo del plan.

% TODO: Incorporar un ejemplo real de respuesta del cliente GPT y el
% serializador utilizado para entregarla al ejecutor. No presentar como
% inferencia observada un plan construido solo para ilustrar el contrato.

\subsection{Tratamiento de aclaraciones y rechazos}

El contrato utiliza \texttt{clarification} cuando falta información
necesaria que puede solicitarse mediante una pregunta breve, y
\texttt{rejected} cuando la tarea excede el dominio o las capacidades
admitidas. En ambos casos, los campos de objetivos y plan se devuelven
vacíos y no se incluyen las marcas de un plan ejecutable.

El consumidor debe utilizar el estado para decidir si existe un plan
susceptible de validación y ejecución. Una pregunta de aclaración no
se interpreta como una acción adicional del plan. La evaluación de
estas respuestas se separa de la generación de acciones para las
instrucciones resolubles del dominio común.

% TODO: Precisar el manejo implementado de estos estados. No afirmar
% que existe un diálogo de aclaración de varios turnos si solo se
% devuelve una pregunta en la respuesta inicial.

\subsection{Validación determinista}

La integración prevé una validación posterior a la generación para
comprobar el estado de respuesta, la estructura de los pasos y su
compatibilidad con el repertorio de acciones. Para los planes
ejecutables se verifican también el orden, los parámetros y las marcas
de inicio y cierre. Las respuestas de aclaración y rechazo deben
mantener vacíos los campos operativos.

La validación de objetivos puede utilizar el contrato semántico común.
La consistencia de las acciones requiere controles específicos sobre
los efectos esperados y las dependencias del plan. Esta comprobación
no delega la planificación a CLIPS ni convierte la arquitectura directa
en una cadena GPT--CLIPS.

Se conservará el resultado original junto con las incidencias del
validador. Las comprobaciones implementadas deberán distinguirse de
las condiciones que solo se solicitan en el prompt, para identificar
qué restricciones se comprueban efectivamente fuera del modelo.

% TODO: Incorporar el validador real de planes y enumerar sus controles.
% El validador de objetivos y las instrucciones del prompt no demuestran
% que exista una validación determinista completa de acciones.

\section{Reproducibilidad}
\label{sec:reproducibilidad_impl}

\subsection{Versiones de modelos}

La identificación de Qwen comprende la versión del modelo base, el
tokenizador y el adaptador seleccionado. El registro deberá permitir
relacionar estos artefactos con el conjunto de datos y la ejecución de
entrenamiento. La conservación de un checkpoint aislado sin su
configuración no permite reconstruir toda la condición experimental.

Para GPT se registrarán el identificador del modelo utilizado en la
solicitud, la fecha de ejecución y las versiones del prompt y del
esquema. La comparación se realizará sin cambiar estas condiciones
durante el procesamiento del conjunto de prueba.

\subsection{Configuraciones experimentales}

La configuración de cada ejecución incluirá los catálogos, los
parámetros de normalización e inferencia, la base de reglas, el modo
de manipulación y el despliegue de los componentes. El mismo criterio
se aplica a los validadores y a los recorridos de conversión de
objetivos.

Para la inferencia se conserva también la distribución de dispositivos,
el número de hilos de CPU y la opción TF32. El módulo configura por
defecto los hilos de cómputo a partir del número de CPU disponibles y
limita los hilos de interoperación a ocho. Cuando CUDA está disponible,
habilita TF32 salvo que se indique \texttt{--no-tf32}. Estas opciones
pueden afectar los tiempos y forman parte de la condición evaluada.

El dataset y las particiones se vincularán con sus tamaños, semillas
y versiones. Las pruebas de funcionamiento se mantendrán identificadas
como parte del desarrollo, mientras que el conjunto de comparación
se gestionará de acuerdo con las condiciones del capítulo 6.

% TODO: Incorporar identificadores o hashes de archivos, commits de los
% componentes y comandos efectivos de lanzamiento, sin credenciales.

\subsection{Registro de ejecuciones}

El sistema simbólico conserva el plan publicado en
\texttt{plan\_execution.log}. Para la evaluación se completará este
registro con el identificador de la orden, la entrada efectiva, los
objetivos originales, las transformaciones del adaptador, el plan y
las incidencias de cada etapa.

El módulo de inferencia incorpora \texttt{ResourceTimer} para registrar
tiempo transcurrido, tiempo de CPU del proceso, memoria RAM máxima del
proceso y memoria de GPU asignada y reservada, junto con sus máximos.
Las mediciones temporales sincronizan CUDA en los límites del bloque
observado. El porcentaje de CPU expresa la razón entre tiempo de CPU y
tiempo transcurrido, por lo que puede superar el 100\% si intervienen
varios hilos. El máximo de RAM corresponde a la vida del proceso, no al
incremento causado únicamente por una solicitud.

Estos diagnósticos se imprimen en la ejecución del módulo y se distinguen
de \texttt{elapsed\_ms}, que el servidor devuelve para la traducción.
El servidor no incorpora automáticamente todas las métricas de
\texttt{ResourceTimer} a cada respuesta HTTP. La latencia completa de
Qwen--CLIPS requiere además medir la comunicación, la conversión de
objetivos y la expansión simbólica con límites temporales explícitos.

En la arquitectura directa se conservarán la respuesta original y el
resultado de las comprobaciones, junto con la latencia y la información
de uso de tokens disponible. Los errores de solicitud y las salidas
inválidas se incluirán en el registro para evitar que la evaluación
considere únicamente los casos que produjeron un plan aceptado.

Esta información permite relacionar un resultado con la configuración
que lo produjo y separar interpretación, conversión, planificación y
ejecución. El capítulo siguiente establece el procedimiento experimental
con el que se evaluarán esas etapas.

% TODO: Confirmar el formato de registro realmente implementado y la
% asociación inequívoca de entrada, salidas y tiempos. El archivo de planes
% existente no equivale por sí solo al registro experimental completo.
